\subsection{The three-division comparison with the signed frames}

\begin{result}{A.3}{Proved}{EXIT\_TOWER E.49, with its label comparisons repaired}
The signed frames used in EXIT\_TOWER E.49 are exactly those of LM7\_PUSH E.65. Its three-division comparison proves equality of their two parities on the multiplier-one subgroup of $\GT$. We supply the actual intermediate maps below. Two label issues in the written argument are repaired: the determinant-minus-one correction is excluded at multiplier one, and the common permutation of the three newly added nonzero two-torsion branch labels is retained until those labels are forgotten.
\end{result}

\paragraph{The common normalization.}
Take $F=(1,f)\in\GT$ and its genuinely cusp-normalized automorphism
$\alpha$ of $\mathcal B=\widehat{C_2(s)*C_3(r)}$, with $\tau=sr$.
In the actual pure subgroup $\widehat F_2=\langle x,y\rangle$,
$x=\tau^{-2}$ and $y=sxs$, the unique pure frames satisfy
\[
 \alpha(\tau)=\tau,\qquad
 \alpha(s)=\beta s\beta^{-1},\qquad
 \alpha(r)=ArA^{-1}.
 \tag{A.C1}
\]
The uniqueness is proved in A.2. With
\[
 U=\begin{pmatrix}1&1\\0&1\end{pmatrix},\quad
 L=\begin{pmatrix}1&0\\-1&1\end{pmatrix},\quad
 S=ULU,\quad R_+=UL,\quad R_0=-R_+,
\]
write $(\beta^{-1})^{\mathrm{ab}}=(\rho,\rho)$ and set
\[
 V=U^{8\rho},\qquad
 T=q_{12}(A)^{-1}V,\qquad W=q_{12}(\beta)^{-1}V.
 \tag{A.C2}
\]
Here $q_{12}$ on a pure word means its signed evaluation at
$x\mapsto U^2$, $y\mapsto L^2$. This is the lift fixed by the
specified modulus-four normalization; we do not discard its overall
sign. The refereed E.65 calculation gives $T=W=I$ modulo four,
$T\bmod3\in C_{\SL_2(\mathbb F_3)}(R_0)$ and
$W\bmod3\in\langle S\rangle$.
The bit $\epsilon_3(T)$ is the sign of its action on
$\ker(R_0-I)$, and $\epsilon_2(W)$ is the exponent modulo two in
the order-four group $\langle S\rangle$. Thus
\[
                  \delta(F)=\epsilon_3(T)+\epsilon_2(W).
 \tag{A.C3}
\]
These equalities identify the character and every sign convention
before any geometric comparison is made.

\paragraph{Constructing the torsion configurations.}
The ordinary universal-curve action is the action proved from the
same-parameter relation and the pants-cell calculations above. It is
an action on full profinite mapping class groups, preserves the named
forgetting homomorphisms, and sends a positive peripheral class to its
$\lambda$th power up to its usual fibre conjugacy. At present
$\lambda=1$.

Let $H=(\mathbb Z/3)^2$ translate an oriented torus and mark its regular
nine-element orbit, with zero distinguished. Its normalizer $N$ has
label-zero subgroup $N_0$, and the actual quotient and forgetting maps
give
\[
 N/H\simeq\Gamma_{1,1},\qquad
 N_0\xrightarrow[\text{forget eight labels}]{\sim}\Gamma_{1,1}.
 \tag{A.C4}
\]
The covering of punctured tori is $F_{10}\triangleleft F_2$, defined
by abelianization modulo three. Every quotient-torus mapping class
lifts, and exactly one deck translate fixes zero. On homology its
linear representative is unchanged by quotienting by $H$ or by
forgetting eight points, proving the stated identification of the
map. Pullback of finite quotients through that forgetting map proves
the full induced topology of $N_0$ and of its finite extension $N$.
The finite-$p$ theorem at $p=3$ therefore identifies the completed
normalizer with $\widehat N$.

We include the determinant and centralizer details in the following
subsection; its proof uses this ordinary universal-curve action and
does not use the CM conclusion.
% Contribution to GT_TOWER_CHECK. Requires amsmath, amssymb, amsthm.
% AI-produced by GPT-6 Astra (OpenAI, ChatGPT Pro mode). No priority claimed.
\subsection{The labelled translation subgroup in the three-division comparison}

The following argument repairs one point in EXIT\_TOWER E.49.  That
proof allows a determinant-minus-one relabelling of the nine translation
labels, and later treats each of the four pair labels as fixed.  The
relabelled action would permute those four pairs.  In the multiplier-one
case under consideration the second pure conjugacy class is excluded by
the peripheral orientation character, as we now prove.  All completions
in this subsection are full profinite completions.

\paragraph{The actual pullback groups.}
Let $T$ be an oriented torus, let
$H=(\mathbb Z/3)^2$ act by translations, and mark its regular orbit
$P=\{p_h:h\in H\}$, with $p_0$ distinguished.  Put
\[
 M=\Gamma(T,P),\qquad
 \widetilde M=\Gamma(T,P\sqcup\{*\})_{\{*\}},
 \qquad \pi:\widetilde M\longrightarrow M.
\]
Here $M$ permits permutations of the nine marks and the subscript on
$\widetilde M$ requires the new mark to be fixed.  The ordinary Birman
sequence has kernel
$\Pi=\pi_1(T\setminus P,*)\cong F_{10}$.
The inverse image $E_H=\pi^{-1}(H)$ is the orbifold fundamental group
of the quotient of $T\setminus P$ by $H$.  Since the action is free,
this is the ordinary fundamental group of the once-punctured quotient
torus.  Consequently the covering gives an exact sequence
\[
 1\longrightarrow F_{10}\longrightarrow E_H\cong F_2
       \xrightarrow{q_H}H\longrightarrow1.                 \tag{CR.1}
\]
The homomorphism $q_H$ is abelianization modulo three followed by an
isomorphism $(\mathbb Z/3)^2\to H$.  Indeed $H$ is abelian of exponent
three and order nine, and $q_H$ is onto.  The profinite Birman sequence
is exact with the full topology on its kernel.  Since $H$ is finite,
the completion of (CR.1) is therefore the actual inverse image of $H$
in the completed universal-curve extension.

Let $F\in\widehat{\mathrm{GT}}$ have multiplier $\lambda=1$.
We use its already constructed action on the ordinary labelled
universal-curve diagram, including the peripheral-power rule.  This is
the tower action used before E.49; no conclusion about the CM
character is used here.  Choose representatives on $M$ and
$\widetilde M$ so that their square with $\pi$ commutes.  To do this,
lift the target inner comparison through the surjective map
$\widehat\pi$ and adjust the representative on
$\widehat{\widetilde M}$ by that lift.  Hence the action restricts to
an isomorphism from $\widehat E_H$ to the inverse image of $F(H)$.

Boggi--Zalesskii Theorem A(i), at $p=3$, realizes $F(H)$ by a discrete
finite subgroup $H'$.  Its conjugator can be chosen pure while allowing
the discrete subgroup $H'$ to vary: subtract a discrete mapping class
with the same nine-label permutation as the original conjugator.
Then $H'$ has exactly the same elementwise label permutations as $H$.
The subgroup $H'$ acts freely on the torus.  In fact a fixed point of
a nonidentity stabilizer is unmarked, since the nine marks form a free
orbit; marking that point produces a nontrivial order-three subgroup
in its universal-curve inverse image.  The original inverse image
$\widehat E_H\cong\widehat F_2$ is torsion-free, and its image under
the tower isomorphism has the same property.  This excludes such a
stabilizer.  Finite realization identifies the resulting free torus
action with a translation action.

Lift the pure conjugator through $\widehat\pi$.  Conjugation by this
lift gives the actual continuous isomorphism
\[
 \theta:\widehat E_H\xrightarrow{\sim}\widehat E_{H'}.
                                                               \tag{CR.2}
\]
Changing the lift changes (CR.2) by an inner automorphism of the
target.  Let $q_{H'}$ identify the deck group with $H$ through its
faithful, unchanged action on the nine labels.  The square on deck
groups then commutes with the identity of $H$:
\[
                       q_{H'}\theta=q_H.                       \tag{CR.3}
\]
This equality follows from preservation of the finite label quotient
and the purity of the correcting conjugator.

\paragraph{The peripheral determinant.}
Write $d_H$ and $d_{H'}$ for positive peripheral loops about the unique
puncture of the two quotient tori.  Every one of the nine upstairs
punctures is a degree-one lift of that puncture.  Its positive loop,
included from the surface kernel into $E_H$, is therefore
$E_H$-conjugate to $d_H$, with exponent one.  The same assertion holds
for $H'$.

The tower action sends each named surface peripheral loop to a
conjugate of its $\lambda$th power in the surface kernel.  Here the
peripheral is the point-pushing loop represented in the universal
curve by the positive twist surrounding the named old mark and the
moving mark.  The labelled universal-curve action preserves that
inertia and its exponent.  Inner comparisons in the base introduce
only the usual surface-peripheral conjugacy: a pure base mapping
class preserves each named peripheral conjugacy class in the fibre.
Thus this is a statement in the fibre group, with its inner ambiguity,
and not merely a conjugacy assertion in the larger mapping class
group.  Applying it to the mark $p_0$ and using the degree-one
inclusion gives
\[
                  \theta(d_H)\sim_{\widehat E_{H'}}
                                     d_{H'}^{\lambda}.           \tag{CR.4}
\]

Choose positive oriented bases $u,v$ and $u',v'$ of the quotient torus
groups.  The same boundary convention identifies the two peripherals
with $[u,v]$ and $[u',v']$, or with their inverses simultaneously.
Let $A\in\operatorname{GL}_2(\widehat{\mathbb Z})$ be the matrix of
$\theta$ on abelianization.  In the quotient by the closed third
lower-central subgroup, expansion of commutators gives
\[
 [u'^{a}v'^{c},u'^{b}v'^{d}]
                   =[u',v']^{ad-bc},\qquad
 \theta([u,v])=[u',v']^{\det A}.                                \tag{CR.5}
\]
The commutator subgroup of this quotient is central, so both
conjugation and the commutator factors in the values of $\theta(u)$
and $\theta(v)$ disappear from this calculation.  The element
$[u',v']$ has infinite profinite order: in
$\operatorname{UT}_3(\mathbb Z/n)$ its image has order $n$, for every
positive integer $n$.  Equations (CR.4)--(CR.5) consequently imply
\[
                             \det A=\lambda=1.                  \tag{CR.6}
\]
This argument uses only the actual peripheral isomorphism (CR.2),
not the CM comparison later in E.49.

Reduce (CR.3) on abelianization modulo three.  The two deck monodromy
isomorphisms satisfy
\[
                  \overline q_{H'}\,(A\bmod3)
                                      =\overline q_H.
                                                               \tag{CR.7}
\]
By (CR.6), their determinants agree.  Oriented quotient-torus basis
changes form $\operatorname{SL}_2(\mathbb Z)$, whose reduction is
onto $\operatorname{SL}_2(\mathbb F_3)$.  The monodromy isomorphisms
form $\operatorname{GL}_2(\mathbb F_3)$, so their two orbits are
exactly their two determinant values.  Equality in (CR.7) therefore
places $H'$ in the same pure conjugacy class as $H$.  The quotient
homeomorphism lifts uniquely after fixing label zero, and that lift
fixes every other label because it intertwines the named deck group.
Thus a pure discrete correcting mapping class suffices.  Together
with the preceding pure profinite conjugator it makes $F(H)=H$
pointwise and preserves all nine labels.

For comparison, the determinant-minus-one coset in
$\operatorname{GL}_2(\mathbb F_3)$ acts by odd permutations on
$\mathbb P^1(\mathbb F_3)$, the four pair labels.  This finite assertion
is checked for all $48$ matrices by the accompanying program.  It
explains why allowing that correction without the determinant
argument does not justify the later assertion that each of the four
forgetting kernels is fixed.

\paragraph{The normalizer topology used in the construction.}
For the translation group, put $N=N_M(H)$ and let $N_0$ fix label zero.
The quotient cover gives
\[
 1\longrightarrow H\longrightarrow N\longrightarrow\Gamma_{1,1}
       \longrightarrow1,
 \qquad N_0\xrightarrow{\sim}\Gamma_{1,1}.
                                                               \tag{CR.8}
\]
Every quotient mapping class preserves the characteristic covering
kernel.  Its lifts form one $H$-torsor, and a unique lift fixes label
zero.  On linear torus representatives the quotient and the map
forgetting the other eight marks induce the same integral homology
matrix.  Thus the second isomorphism in (CR.8) is the actual ordinary
forgetting map.  Every finite quotient of $N_0$ pulls back through
that ambient map.  Hence $N_0$ has the full induced topology.
Its finite-index inclusion in $N$ is induced-open because it
stabilizes a finite label, so $N$ has full induced topology as well.
Boggi--Zalesskii Theorem A(iii), now at $p=3$, gives
\[
                        N_{\widehat M}(H)=\widehat N.           \tag{CR.9}
\]
The genus-one mixed mapping class group is good and of finite virtual
cohomological type: start with the virtually free group
$\Gamma_{1,1}=\operatorname{SL}_2(\mathbb Z)$ and use the full
profinite Birman sequences with finitely generated free kernels.
Filtered unions of finite modules give the arbitrary discrete
continuous module condition in Theorem A, as in the genus-zero
centralizer argument.

\paragraph{Successive fixed-point centralizers.}
The center of $N_0$ contains its unique nonidentity elliptic
involution $\iota$.  The nine original marks consist of one fixed
mark and four pairs.  On adjoining the other fixed points, the
successive orbifold inverse-image groups are
\[
 (C_2*C_2*C_2)*F_4,\quad (C_2*C_2)*F_5,\quad C_2*F_6,\quad F_7.
                                                               \tag{CR.10}
\]
Their even kernels have ranks $10,11,12,13$.  The free-product torsion
calculation, followed by Theorem A at $p=2$, gives respectively
three, two, one, and zero remaining fibre-conjugacy classes of
order-two lifts.  A chosen lift $x$ has centralizer $\langle x\rangle$
in the orbifold inverse image; its intersection with the even kernel
is trivial.

We record the projection comparison explicitly.  In an extension
$1\to K\to E\to L\to1$, the projection of $C_E(x)$ is the stabilizer
of the $K$-conjugacy class of $x$.  An element of that stabilizer has
a lift $e$ with $exe^{-1}=kxk^{-1}$ for some $k\in K$, and
$k^{-1}e$ centralizes $x$.  The kernel is $C_K(x)$, which is trivial
in (CR.10).  Thus the centralizer maps isomorphically to the relevant
finite-index class stabilizer.  Every finite quotient pulls back
through ordinary forgetting to the preimage of that stabilizer in
the ambient group.  This proves its full induced topology and
identifies its completion with the actual completed pullback.
Taking ordered choices and then their finite permutation descent
therefore adjoins the three fixed points functorially.

\paragraph{The final involution quotient.}
After adjoining these points, let $G$ preserve the four fixed labels
individually and preserve the four named pairs, allowing each pair to
flip.  For $C=C_G(\iota)$ the actual unramified double cover gives
\[
 P_0=F_{13}\mathrel{\triangleleft}E_0=F_7,\quad E_0/P_0=C_2,
 \qquad
 1\longrightarrow\langle\iota\rangle\longrightarrow C
       \longrightarrow\Gamma_{0,8}\longrightarrow1.            \tag{CR.11}
\]
At branch cusps the positive upstairs peripheral is the square of
the positive downstairs peripheral; unique roots in the discrete
free group $E_0$ justify descent.  At pair cusps the local degree is
one.  The pure downstairs group preserves the parity homomorphism,
so every such mapping class lifts; the kernel consists of deck
transformations.

Write $P=\widehat P_0$ and $E=\widehat E_0$.  Their extension has
$C_E(P)=1$: a centralizing element has square in $Z(P)=1$, while
$E$ is torsion-free.  Consequently $E$ is the inverse image of the
deck subgroup in $\operatorname{Aut}(P)$.  A centralizing outer
automorphism of $P$ therefore descends to $\operatorname{Out}(E)$;
changing its representative changes the descendant only by an inner
automorphism.  If a class from $\widehat C$ acts trivially on $P$,
adjust its representative by an element of $P$ to make that action
the identity.  Its action on $E$ is then the identity because
$C_E(P)=1$.  Genus-zero congruence faithfulness places the class in
the finite deck kernel of the completion of (CR.11).  That kernel
is detected by the action $-I$ on torus homology modulo three.
Thus $\widehat C$ embeds in the actual ambient $\widehat G$ and has
its full induced topology there.  Direct Theorem A(iii), for the
single order-two subgroup, proves
\[
 C_{\widehat G}(\iota)=\widehat C,\qquad
 \widehat C/\langle\iota\rangle\cong\widehat\Gamma_{0,8}.        \tag{CR.12}
\]
There is no exchange of intersection and closure in this argument.

Each of the four ordinary downstairs marks can be moved once around
a branch mark; its two lifts exchange and the other three pairs stay
fixed.  These four independent monodromies connect all individual
sheet-order choices.  Hence the involution with the specified four
fixed labels and four pairs has one pure conjugacy class.  It can be
straightened by a pure conjugator.

The four maps used by E.49 are now individually named: for
$i=1,2,3,4$, retain the four $E[2]$ branch labels and the $i$th
$E[3]/\{\pm1\}$ pair label and forget the other three pairs.
The determinant repair pins those four pair labels.  The ordered
construction of the three additional order-two lift classes still
permits one common permutation of their names.  We retain this
permutation instead of asserting that an arbitrary correcting
inner automorphism of the original $\Gamma_{1,1}$ fixes their
individual names.

Write the induced eight-point outer action as $(F',\sigma)$ in
the genus-zero decomposition
$\widehat{\mathrm{GT}}\times\mathfrak S_8$.
The four forgetting kernels make $\sigma$ preserve each retained
five-element set, hence fix the four pair labels and preserve the
branch quartet.  The original branch label zero is fixed by the
construction, so the remaining permutation lies in the
$\mathfrak S_3$ of the three newly added branch marks.
On each five-point target E.44 identifies the GT parameter with
$F$; geometric relabelling changes the permutation factor and
does not change that GT factor.  Thus $F'=F$ and the exact
eight-point conclusion is $(F,\sigma)$ with this one common
$\sigma\in\mathfrak S_3$.

The map used subsequently in E.49 retains the original branch
label zero and the four pair labels, and forgets precisely the
other three branch marks.  The preceding $\sigma$ disappears
under this map.  Explicitly, choose a disk containing those three
marks and none of the five retained marks, and represent $\sigma$
by a braid supported in that disk.  After forgetting the three
marks its mapping class is the identity.  Hence the induced
five-point outer action is exactly $(F,1)$, which is the parameter
and labelling required for the monic division-quartic construction.
Equivalently one may remove $\sigma$ at the eight-point stage by
this one geometric relabelling, keeping its induced source inner
comparison.  The proof does not require that all twelve labels
were already pinned by the original translation correction.

Finally this correction is compatible with the single global gauge
in E.267.  After rigidifying the original central involution, write
the corrected source representative as
$\alpha'=\operatorname{Ad}(\gamma)\alpha$ with
$\gamma\in\widehat B$.  For the actual map
$\varphi:\widehat B\xrightarrow{\sim}G_3\subset G_2$,
the inner comparison has the exact algebraic transformation
\[
 \operatorname{Ad}(h)\varphi\alpha'
   =\operatorname{Ad}\bigl(h\varphi(\gamma)\bigr)
       \varphi\alpha.                                      \tag{CR.13}
\]
The multiplier-one source action is the identity on the
abelian quotient $B^{\rm ab}=C_6$, and
$q\varphi(\gamma)$ belongs to the same actual $C_6$ subgroup.
Thus the adjusted comparison still supplies one exponent in
$C_6$.  In particular no independent local signs are introduced
by the remaining branch permutation or by the source inner
ambiguity.


\paragraph{The four comparisons, with the residual permutation retained.}
After adjoining the other three fixed points of the elliptic
involution, the twelve marks are $E[2]\cup E[3]$. Their involution
quotient has four branch labels $b_0,b_1,b_2,b_3$, with $b_0$ the
original origin, and four pair labels $t_j$, indexed by
$j\in\mathbb P^1(\mathbb F_3)$. The preceding construction fixes
each original three-torsion label. It fixes $b_0$; the other three
branch labels may have one common geometric permutation.

For each $j$ the comparison homomorphism is exactly
\[
 p_j:\widehat\Gamma_{0,8}\longrightarrow
 \widehat\Gamma_{0,\{b_0,b_1,b_2,b_3,t_j\}},
 \tag{A.C5}
\]
obtained by forgetting the other three pair labels. These are the
four uses of the branch-order comparison of A.1. In each, the
selected pair is the two-point elliptic input, the four fixed points
give its branch-order cover, and the five-point GT parameter is
exactly $F$. A different ordering of the three new fixed points
inserts only its stated geometric permutation factor. It changes
no GT parameter.

For completeness, (P.1) writes the outer action on the eight-point
group as $(F',\sigma)\in\GT\times\mathfrak S_8$. Each of the four
kernels in (A.C5) is retained: a two-point twist survives the
forgetting map precisely when both of its labels survive. Thus
$\sigma$ preserves the four retained five-element sets. It fixes
all four $t_j$ and preserves the branch quartet. Since the origin
was individually retained throughout, it also fixes $b_0$.
Restriction to any one quotient and its comparison in A.1 gives
$F'=F$. The only remaining permutation is therefore in
$\mathfrak S_{\{b_1,b_2,b_3\}}$.

Now use the actual forgetting homomorphism that removes
$b_1,b_2,b_3$ and retains $b_0,t_1,t_2,t_3,t_4$. A permutation of
the three forgotten labels can be represented by braids in a disk
containing those three and none of the retained points. This
representative becomes the identity after forgetting them. Its
arbitrary other representatives contribute only an inner
automorphism of the retained group. Consequently the resulting
five-point action has parameter $F$ and the required distinguished
origin. This is the comparison needed for the quartic below. It
avoids an unsupported assertion that all the newly chosen branch
labels were individually fixed from the start.

\paragraph{The division quartic and its ordered poles.}
For the elliptic curve $y^2=x^3+ax+b$, let $D=4a^3+27b^2\ne0$.
The four images of its nonzero three-torsion points are the roots
of $3x^4+6ax^2+12bx-a^2$. The quartic double cover has the explicit
birational change of variables
\[
\begin{gathered}
 t^2=x^4+2ax^2+4bx-a^2/3,\qquad
 X=(t+x^2+a)/2,\quad Z=X-a/3,\quad Y=Xx+b/2,\\
 Y^2=Z^3+D/108,\qquad P=(-a/3,-b/2),\\
 x=\frac{Y-b/2}{Z+a/3},\qquad t=2(Z+a/3)-x^2-a.
\end{gathered}
\tag{A.C6}
\]
Substitution proves both inverse identities in the function fields.
They therefore extend to inverse morphisms of the smooth projective
curves. The positive pole of the monic quartic is the origin $O$,
and the negative pole is $P$. Its deck involution is $Q\mapsto P-Q$:
the divisor of the degree-two function $x$ has poles $O+P$, so
points in one fibre have group sum $P$. Thus its branch points
satisfy $2Q=P$. At $b=0$ the point $P$ is nonzero two-torsion.
At $a=0$ its horizontal tangent has intersection multiplicity three,
so it is nonzero three-torsion.

An original Weierstrass scaling by $u$ induces
$(Z,Y)\mapsto(u^4Z,u^6Y)$. In particular $u=i$ acts by the elliptic
involution and $u=\zeta_6$ acts by multiplication by $\zeta_3^2$.
These calculations fix the two CM signs in this model. The symbolic
forward and inverse checks are in \texttt{check\_cm\_comparisons.py}.

The pole order is intrinsic. Under an affine change $x'=cx+d$, the
monic equation transforms by $t'=c^2t$, and therefore preserves the
sign of $t/x^2$ at infinity. Every permutation of the four roots
preserves the same monic polynomial and these two poles separately.
The target group is consequently the mapping class group with two
ordered elliptic points.

Let $G_2=\widehat\Gamma_{1,2}$, $G_1=\widehat\Gamma_{1,1}$ and
$q:G_2\twoheadrightarrow G_1$ be ordinary forgetting. On the level
$D=108$, the substitutions $a=-3Z$, $b=-2Y$ identify the
discriminant complement with the affine curve
$\mathcal E_0\setminus\{O\}$, where $\mathcal E_0:Y^2=Z^3+1$.
After rigidifying the original central involution, the residual
scaling variable is $v=u^2$, of weights $(2,3)$ on $(a,b)$ and
weight six on $D$. We obtain an actual stack isomorphism and its
fundamental-group homomorphism
\[
 \left[\{D\ne0\}/\mathbb G_m{}_{(2,3)}\right]
   \simeq[(\mathcal E_0\setminus\{O\})/\mu_6],
 \qquad
 \varphi:\mathcal B\xrightarrow{\sim}
       G_3:=q^{-1}(\langle R_+\rangle).
 \tag{A.C7}
\]
On the left, $\mathcal B$ is the full profinite orbifold group of
the rigidified once-pointed elliptic moduli stack. On the right,
the group is the actual pullback of the ordinary universal curve
over its CM stabilizer.

The topology in (A.C7) is full: the discrete pullback has free kernel
$F_2$ and quotient $C_6$. That free kernel has full topology because
it has finite index in the source. On the target side, the profinite
Birman sequence embeds the same $\widehat F_2$ as $\ker q$, and the
CM subgroup embeds as its finite quotient. The induced completed
map is the identity on the kernel and an isomorphism on the quotient,
so is an isomorphism. This argument does not treat $G_3$ as an open
subgroup of $G_2$. The commutator subgroup $K_0=\mathcal B'$ is
identified with $K=\ker q$. On the cyclic quotient,
$\varphi(\tau)$ maps to $R_+^{\varepsilon}$, with the actual
$\varepsilon\in\{1,-1\}$ retained below.

Every comparison leading to (A.C7) has now been specified:
translation normalizer and ordinary forgetting, fixed-point lift
classes, the completed involution centralizer quotient, the four
maps (A.C5), forgetting the three extra branch values, and the
ordered-pole double cover. Composing their outer-commutative squares
leaves one inner comparison in the final group $G_2$. Inner changes
on the source pass through $\varphi$ and are absorbed into that
same final comparison. This is the naturality used below.

\paragraph{The point-pushing matrix is the stated signed shear.}
Use the braid presentation
\[
 G_2=\widehat B_4/\overline{\langle\omega_4\rangle},\qquad
 G_1=\widehat B_3/\overline{\langle\omega_3^2\rangle},\qquad
 (a,b,c)\longmapsto(u,w,u).
\]
The kernel $K$ is free on
$\xi_1=ca^{-1}$ and $\xi_2=aba c^{-1}b^{-1}a^{-1}$. Its exact
conjugation table is
\[
\begin{array}{c|cc}
 &\xi_1&\xi_2\\ \hline
 a&\xi_1&\xi_2\xi_1\\
 b&\xi_1\xi_2^{-1}&\xi_2\\
 c&\xi_1&\xi_1\xi_2
\end{array}
\tag{A.C8}
\]
Thus $a,c$ act by $U$ and $b$ by $L$ on $K^{\mathrm{ab}}$.
The table follows by substitution in the braid presentation and
also reconstructs $B_4=F(\xi_1,\xi_2)\rtimes B_3$. The central
word $\omega_4$ maps to the infinite-order $\omega_3^2$ before the
quotients, so quotienting retains the free kernel. The full
completed comparison is the ordinary Birman sequence already used.
The exact reduced-word checks are retained with the braid programs.

For the genuine four-strand action $\Phi$, put
$e=f(ba^2b,c^2)$. Since $ba^2b=(ab)^3a^{-2}$, $a$ commutes with
both arguments. Removing its commuting central factor from a
commutator proword gives $e=f((ab)^3,c^2)$ and $[e,a]=1$.
After filling, these arguments commute; hence $e\in K$.
The formulas $\Phi(a)=a$ and $\Phi(c)=e^{-1}ce$ imply, for the
matrix $N$ on $K^{\mathrm{ab}}$,
\[
 (U-I)[e]=0,\qquad Ne_1=e_1,\qquad NU=UN.
\]
Writing out a two-by-two matrix gives $N=U^t$ for one
$t\in\widehat{\mathbb Z}$. Covariance with $b$ gives
\[
 NLN^{-1}=q_{12}(f)^{-1}Lq_{12}(f)\pmod {12}.
\]
The signed identity in the refereed E.65 gives the same equation
for $V=U^{8\rho}$, and the independent finite enumeration checks
it for every permitted signed residue. The common centralizer of
$U,L$ consists of scalar matrices: commuting with $U$ first makes
the lower-left entry zero and the diagonal entries equal; commuting
with $L$ then kills the upper-right entry. Since both $N$ and $V$
fix $e_1$, their scalar quotient is the identity. Therefore
\[
                           N\equiv V\pmod {12}.
 \tag{A.C9}
\]

Lift the pure frame $A$ to $\widetilde A\in G_1$ by its actual
word in $u^2,w^2$, and lift that same word to $G_2$ using $a^2,b^2$.
Its signed matrix is $q_{12}(A)$. The projected frame equation gives
$\Phi(R_+)=\widetilde A R_+\widetilde A^{-1}$ up to the central
involution. That remaining sign is fixed by
$G_1^{\mathrm{ab}}=\mathbb Z/12$: the multiplier-one action is
the identity there, whereas multiplying by the central involution
adds six. Consequently
\[
 \Phi^A=\operatorname{Ad}(\widetilde A^{-1})\Phi
 \quad\hbox{preserves }G_3,\qquad
 (\Phi^A)_{K^{\mathrm{ab}}}\equiv T\pmod {12}.
 \tag{A.C10}
\]
This proves the exact occurrence of the frame $T$ in the geometric
comparison.

\paragraph{The commutator homology of the source.}
The source has the explicit metabelian quotient
\[
 \begin{gathered}
 \mathcal B/[K_0,K_0]
   =\widehat{\mathbb Z}^{\,2}\rtimes C_6,\qquad
 R=\begin{pmatrix}0&-1\\1&1\end{pmatrix},\\
 \tau=(0,R),\quad s=(p,R^3),\quad r=(p,R^4),\quad p=e_1.
 \end{gathered}
 \tag{A.C11}
\]
One can verify this without assuming a split extension: in the
Schreier presentation put $p_i=\tau^i(s\tau^{-3})\tau^{-i}$ and
$z=\tau^6$. The relations become
$p_i+p_{i+3}+z=0$ and $p_i+p_{i+4}+p_{i+8}+2z=0$. They give
$z=0$ and $p_2=p_1-p_0$. The displayed semidirect product proves
that $p_0,p_1$ are independent and gives $R$ as claimed. The
index-six torsion-free subgroup is free of rank two; passing to
its full completion gives (A.C11).

Let $M$ be the matrix induced by $\alpha$ on $K_0^{\mathrm{ab}}$.
Since $\alpha(\tau)=\tau$, it commutes with $R$. Write
$\beta=(b,R^j)$ and $A=(a,R^\ell)$ in (A.C11). Purity makes
$j,\ell$ even. The frame equations give
\[
 Mp=2b+R^jp=(I-R^4)a+R^\ell p,\qquad j=4\rho\pmod6.
 \tag{A.C12}
\]
The last equality is the abelian exponent calculation: both $x$
and $y$ have exponent $-2$ in the cyclic quotient, and $\beta$
has exponent vector $(-\rho,-\rho)$.

Modulo three, $R^4=-R$, and $R+I$ has rank one, is nilpotent,
and has equal image and kernel $L_-$. The second expression in
(A.C12) shows that $M$ is $+1$ on the quotient by this line,
because $\ell$ is even. Commutation with $R+I$ gives the same
scalar on the line.

For the modulus-four calculation, identify the projective action
$\PSL_2(\mathbb F_3)=A_4$ with its affine action on
$\mathbb F_2^2$ by
$\infty\mapsto0$, $0\mapsto p$, $-1\mapsto Rp$,
$1\mapsto p+Rp$. Then $S,R_0,U^{-1}$ act respectively as
$(p,I),(p,R^4),(0,R)$. The projective image of $W$ is translation
by $-R^{-j}b$, and its membership in $\langle S\rangle$ says
$b=\epsilon_2(W)R^jp$ modulo two. Substitution in the first
expression of (A.C12), followed by commutation with $R$ on the
basis $p,Rp$, gives
\[
 M|_{L_-}=1\pmod3,\qquad
 M\equiv(-1)^{\epsilon_2(W)}R^j\pmod4,\qquad j\equiv0\pmod2.
 \tag{A.C13}
\]
These formulas retain the full source automorphism; they are
necessary images under the explicit metabelian and coefficient
quotient maps, rather than a lifting assertion from those quotients.

\paragraph{One conjugator gives one exponent at both primes.}
The composed naturality comparison explained after (A.C7), now
with the correction (A.C10), supplies a single $h\in G_2$ such
that
\[
                  \Phi^A\varphi
                   =\operatorname{Ad}(h)\varphi\alpha.
 \tag{A.C14}
\]
Both sides induce the identity on the cyclic CM quotient
$\langle R_+\rangle$. Thus $q(h)$ centralizes $R_+$ in $G_1$.
Modulo the central involution, this is the centralizer of the
order-three free factor in $\widehat{C_2*C_3}$; the finite-$p$
theorem at $p=3$ identifies it with that factor. Its full preimage
in $G_1$ is $\langle R_+\rangle$. We therefore have
\[
                 h\in G_3,\qquad q(h)=R_+^k,\qquad k\in\mathbb Z/6.
 \tag{A.C15}
\]
There is one exponent $k$ because there is one group element $h$.
Taking separate conjugators after reduction at two and three would
not give (A.C15), and is not used.

The actual discrete map behind $\varphi$ induces an integral
isomorphism $C:K_0^{\mathrm{ab}}\to K^{\mathrm{ab}}$, with
$CR=R_+^{\varepsilon}C$. Applying abelianization and reduction
modulo twelve to (A.C14) gives
\[
                         TC=R_+^kCM\pmod {12}.
 \tag{A.C16}
\]
At three, $C$ sends $L_-$ to $\ker(R_0-I)$. Formula (A.C13)
makes $M$ act by $+1$ on that line and $R_+$ acts by $-1$.
Hence $\epsilon_3(T)=k$ modulo two. At four, $T=I$ and
$R^3=-I$. Formula (A.C13) has exponent $j+3\epsilon_2(W)$ in
the cyclic group $\langle R\bmod4\rangle$ of order six.
Equation (A.C16) gives $\epsilon_2(W)=k$ modulo two.
The possible $\varepsilon=-1$ changes neither parity. We obtain
\[
                   \epsilon_3(T)=\epsilon_2(W),\qquad\delta(F)=0.
 \tag{A.C17}
\]

\paragraph{The pentagon line and the witness.}
The ordinary tower construction uses the genus-zero GT action,
including its coherent action on the ten-point configuration in
A.1. In the explicit local proof, the pentagon is used at the
line reducing the six-move word to
\[
 f(W,X)f(Y,Z)f(V,W)f(X,Y)f(Z,V)=1
\]
by the five-point pentagon, the line corresponding to EXIT\_TOWER
(E.98). That is the extra relation used to construct the ordinary
sewing-and-forgetting action. In the final CM comparison its exact
consequence is the single naturality equation (A.C14).

For every witness in A.2, $\rho=0$, $T=5I$, $W=I$.
Formula (A.C13) then gives $M=I$ modulo four and $M|_{L_-}=1$
modulo three. Equation (A.C16) would require $k$ even at four and
odd at three. The enumerated signed matrix certificate verifies
the absence of such a $k$. Therefore (A.C14) fails for that
witness. The already proved implication from the full pentagon
to (A.C14) excludes its full pentagon. This locates the failed
naturality equation and the pentagon line that supplies it; it
does not pretend to have computed a particular finite five-point
pentagon value for the existential word $f$.

There is also an exact identification of the group in that line.
The generators $X,Y,Z,V,W$ lie in the pure subgroup
$P_4/\langle\omega_4\rangle=\Gamma_{0,5}$ of
$B_4/\langle\omega_4\rangle=\Gamma_{1,2}$. Its index is
$24$, so completing the inclusion is injective and has open
image. The substitutions $V=Xx_{13}Y$ and $W=Yx_{24}Z$
identify the displayed word with the standard pentagon under
this inclusion. Consequently the nontrivial full pentagon of
the witness remains nontrivial in the very group of (E.98).
This proves failure at the stated line; it uses the proved
comparison and the finite-index injection, without assigning
an uncomputed finite quotient to that nontrivial word.

\begin{result}{A.7}{Proved}{The CM character vanishes on the full profinite GT group}
With the character of the refereed LM7\_PUSH E.65 and its stated
signed normalization,
\[
 \boxed{\ \Delta|_{\GT}=0\ },\qquad
 \GQ\subseteq\GT\subsetneq\ker\Delta\subsetneq\GTz.
 \tag{A.C18}
\]
The last two inclusions are strict. The first is the
Drinfeld--Ihara embedding; its surjectivity is not decided
in this audit.
\end{result}
\begin{proof}
Equation (A.C17) proves the multiplier-one assertion. For an
arbitrary $F\in\GT$, surjectivity of the cyclotomic character
provides $\sigma\in\GQ$ with $\lambda(\sigma)=\lambda(F)$.
The actual product $F\sigma^{-1}$ lies in $\GT$ and has
multiplier one. The supplied character is a homomorphism, is
zero on $\GQ$, and restricts to $\delta$ at multiplier one.
Thus $\Delta(F)=\delta(F\sigma^{-1})+\Delta(\sigma)=0$.
The standard Galois embedding gives the first inclusion.
The second containment is the proved vanishing, and B.2.1
supplies an actual full profinite gentle element in its
complement inside $\ker\Delta$. Finally A.2 gives an element
of $\GTz$ with character one, proving the last strict inclusion.
\end{proof}
